% Paper A: conclusions.
\section{Conclusions}\label{sec:conclusion}
The factorization of the symbol of a linear constant-coefficient PDE is a grouping structure on its solutions: the
modes of one factor lie on one curve of the dispersion relation, the solution splits into one component per factor up
to an overlap of dimension $\deltaP$, and the components are as well separated as the Friedrichs angle between their
solution spaces allows. Additive terms do not form groups, and their interaction cannot be regularized away. These
facts determine what a learning algorithm should look for and what it can promise; \cref{tab:status} separates what is
proved from the design rules derived from it and from what is measured.

% Shared: the kinds of claim (proved / design rule / measured).
\begin{table}[tbp]
\centering\small
\caption{The kinds of claim made in the paper.}\label{tab:status}
\begin{tabularx}{\linewidth}{@{}>{\raggedright\arraybackslash}p{2.2cm}X@{}}
\toprule
Kind & Statements\\
\midrule
proved & groups are factors with overlap $\deltaP$ (\cref{thm:factor}); terms are not groups (\cref{prop:terms,prop:discont});
identification criterion and its limits on finite data (\cref{thm:ident,prop:unique,prop:certify}); stability via
$\gN$ (\cref{thm:hm}); Bayesian error bars via the angle (\cref{prop:bayes-angle}); counting condition (\cref{prop:counting})\\
design rule & labels instead of ``identified''; distributions of $\rho$ instead of values; decisions on the raw observations;
growth of the class instead of coefficient penalties; minimum-norm split of the overlap\\
measured & $0/20$ confident wrong answers on random equations not used in development; $60/60$ factor sets on grid
designs with independent fields and noise; $30/30$ on the development benchmark; coefficient coverage $0.94$ at
nominal $0.90$\\
\bottomrule
\end{tabularx}
\end{table}


The algorithm built on them groups twice, modes into factors within a sample and factors across samples, and reports a
distribution over equations whose labels and intervals reflect the limits proved in \cref{sec:analysis}: an equation
is never certified, closely aligned groups come with wide intervals, and repeated factors are reported as undecided.
On twenty random equations not used in development it is never wrong without warning; on the thirty-input development
benchmark it recovers every factor set,
where the best regression on a library of terms fails without warning on five inputs, four of them products of
factors.

Three directions remain open. For variable coefficients the eigenvalues of the pencil become functions of position
(\cref{rem:varcoef}), and the factors become families of curves; the algebra is then non-commutative, and a splitting
theorem is not available. For nonlinear equations the ring is replaced by a dictionary and grouping must be relative
to it (\cref{sec:nonlinear}). For systems such as the Dirac equation the ideal becomes a module and the conditions for
splitting become topological. In each case the question asked in this paper, which pieces of the equation a solution
carries and how stably, is the right starting point.
